\section{Síntesis y ampliación}

El hilo conductor de esta entrevista puede resumirse en tres etapas: «primero la plataforma, luego el cambio de paradigma y después la implementación del sistema». La narrativa de Jensen no consiste en predecir alguna tendencia de corto plazo, sino en explicar por qué NVIDIA ha logrado mantenerse en una posición central a lo largo de varios ciclos: construir la plataforma con anticipación y esperar a que el ecosistema amplifique la innovación.

\begin{longtable}{p{0.22\textwidth} p{0.34\textwidth} p{0.34\textwidth}}
\toprule
\textbf{Tema} & \textbf{Postura en la entrevista} & \textbf{Implicación de ingeniería} \\
\midrule
Origen de la computación paralela & Las GPU se escalaron gracias al mercado de los videojuegos; CUDA abrió la programación de propósito general & La inversión en plataformas debe vincularse a una fuente de demanda sostenible \\
\midrule
El momento de AlexNet & Validó Software 2.0 y la ruta del entrenamiento a gran escala & Cuando ocurre un cambio de paradigma, hay que reconstruir toda la pila en lugar de optimizar partes aisladas \\
\midrule
Physical AI & Omniverse + Cosmos sostienen el entrenamiento y la generalización de los robots & El world model y el motor de simulación deben diseñarse de manera conjunta \\
\midrule
Seguridad de la IA & El riesgo se descompone en capa del modelo, capa del sistema y capa de uso & La gobernanza de la seguridad debe ser ingenieril, procesal y auditable \\
\midrule
Restricción energética & El aumento del cómputo está limitado por el consumo eléctrico y la eficiencia del sistema & El foco de la competencia se traslada a la ingeniería de sistemas y a la relación de eficiencia energética \\
\midrule
Evolución de la arquitectura & El transformer no será el punto final & Mantener la flexibilidad de la plataforma es preferible a apostar por una sola arquitectura \\
\midrule
Desarrollo personal & Preguntarse continuamente «¿cómo hacer las cosas mejor con IA?» & La capacidad de usar prompt y flujos de trabajo se volverá una competencia profesional general \\
\bottomrule
\end{longtable}

\begin{importantbox}{Síntesis en una frase}
Si la palabra clave de la última década fue «entrenar modelos grandes», la de la próxima década será, con más probabilidad, «convertir la capacidad de los modelos en un sistema real que sea escalable, gobernable y asequible».
\end{importantbox}

\subsection{Lecturas adicionales}
\begin{itemize}
    \item Documentación oficial de NVIDIA CUDA: \url{https://docs.nvidia.com/cuda/}
    \item Alex Krizhevsky, Ilya Sutskever, Geoffrey Hinton (2012), \textit{ImageNet Classification with Deep Convolutional Neural Networks}
    \item NVIDIA Omniverse: \url{https://www.nvidia.com/en-us/omniverse/}
    \item Material de NVIDIA Robotics/Isaac: \url{https://developer.nvidia.com/isaac}
    \item Hoja de ruta de la arquitectura de GPU de NVIDIA (material público de Hopper/Blackwell): \url{https://www.nvidia.com/en-us/data-center/}
    \item Video de esta entrevista (Cleo Abram): \url{https://www.youtube.com/watch?v=7ARBJQn6QkM}
\end{itemize}

%% --- Fin del contenido --- %%

\end{document}
